\chapter{Endoscopic Ultrasound (EUS)}

\section*{Overview}
Endoscopic ultrasound (EUS) combines endoscopy and ultrasound to obtain detailed images and information about the digestive tract and the surrounding tissue, organs, and lymph nodes.

\section*{Alternative Names}
EUS, Echo-endoscopy

\section*{Principle}
An endoscope with a small ultrasound transducer at its tip is passed into the esophagus, stomach, or duodenum, placing the probe close to the organs. Sound waves create detailed images of the gut wall and adjacent structures such as the pancreas, bile ducts, and lymph nodes, and fine-needle aspiration can sample suspicious lesions.
\section*{History}
Developed in the early 1980s, EUS integrated high-frequency ultrasound transducers onto flexible endoscope tips, enabling high-resolution mediastinal and abdominal imaging.

\section*{Why This Test or Procedure Is Ordered}
Ordered to evaluate pancreatic cysts, tumors, gallstones in the bile duct, determine the depth of gastrointestinal tract cancer invasion, or obtain biopsies of deep tissues.

\section*{Is This a Primary or Secondary Test or Procedure}
Primary diagnostic and staging tool for pancreatico-biliary cancers and submucosal lesions.

\section*{How This Test or Procedure Is Performed}
Performed under deep sedation. A specialized endoscope with an ultrasound probe at the tip is passed down your esophagus or into your rectum to scan adjacent structures.

\section*{What Preparation Is Needed}
Fast for 6 to 8 hours. Arrange for transportation home due to sedation.

\section*{Contraindications and Precautions}
Gastrointestinal perforation, severe coagulation disorders (if fine needle biopsy is planned).

\section*{Risks}
Perforation, bleeding, pancreatitis (especially if biopsy of the pancreas is performed), or sedation complications.

\section*{What Happens After the Test or Procedure}
Monitor in recovery. Avoid driving for 24 hours. Sore throat or mild abdominal discomfort may occur.

\section*{Understanding Your Results}
Provides highly detailed ultrasound images of the walls of the esophagus, stomach, duodenum, pancreas, and bile ducts. Helps stage tumors accurately.

\section*{Normal Results}
Normal wall layers of the GI tract; normal pancreas, bile ducts, and adjacent lymph nodes; no masses or cysts.

\section*{Follow-up Tests and Procedures}
Biopsy review, oncology staging discussion, or surgical consult.

\section*{References}
\begin{enumerate}
  \item MedlinePlus. Endoscopic Ultrasound (EUS). U.S. National Library of Medicine; 2025.
  \item Current Medical Diagnosis and Treatment. McGraw-Hill; 2025.
\end{enumerate}

\clearpage
